%\documentclass[12pt,oneside]{amsart}
\documentclass{scrartcl}

%\documentclass{article}

\usepackage{amsthm,amsmath,amssymb}
\usepackage[numbers]{natbib}

%\usepackage[colorlinks]{hyperref}
\usepackage{hyperref}
\hypersetup{
    colorlinks,%
    citecolor=black,%
    filecolor=black,%
    linkcolor=black,%
    urlcolor=black
}
\usepackage{hypernat}
\usepackage[top=2.5cm, bottom=2.5cm, left=2.8cm,
right=2.8cm]{geometry}
\usepackage{graphicx}


%\setlength{\parskip}{0pt}
%\setlength{\parsep}{0pt}
%\setlength{\headsep}{0pt}
%\setlength{\topskip}{0pt}
%\setlength{\topmargin}{0pt}
%\setlength{\topsep}{0pt}
%\setlength{\partopsep}{0pt}

\linespread{1}

%\usepackage{marvosym}

%\textwidth = 6.5 in \textheight = 9.2 in \oddsidemargin = 0.0 in
%\evensidemargin = 0.0 in \topmargin = -0.0 in \headheight = 0.0 in
%\headsep = 0.0in \parskip = 0.0in \parindent = 0.0in
%\def\baselinestretch{0.871}


\newtheorem{theorem}{Theorem}[section]
\newtheorem{proposition}[theorem]{Proposition}
\newtheorem{problem}[theorem]{Problem}
\newtheorem{fact}[theorem]{Fact}
\newtheorem{lemma}[theorem]{Lemma}
\newtheorem{defn}[theorem]{Definition}
\newtheorem{corollary}[theorem]{Corollary}
\newtheorem{remark}{Remark}[section]
\newtheorem{example}[theorem]{Example}

\newcommand{\E}{{\mathbb E}}
\newcommand{\se}{{\mathcal{E}}}
\newcommand{\R}{{\mathbb R}}
\renewcommand{\P}{{\mathbb P}}
\newcommand{\ac}{{\mathcal{A}}}
\newcommand{\A}{{\mathcal{A}}}
\newcommand{\B}{{\mathcal{B}}}
\newcommand{\C}{{\mathcal{C}}}
\newcommand{\M}{{\mathcal{M}}}
\newcommand{\Mf}{{\mathfrak{M}}}
\newcommand{\T}{{\mathcal{T}}}
\newcommand{\Z}{{\mathcal{Z}}}
\newcommand{\K}{{\mathcal{K}}}
\newcommand{\lc}{{\mathcal{L}}}
\newcommand{\Ps}{{\mathcal{P}}}
\newcommand{\G}{{\mathcal{G}}}
\newcommand{\pa}{{\mathring{p}}}
\newcommand{\F}{{\cal F}}
\newcommand{\samp}{{\mathcal{X}}}
\newcommand{\X}{{\mathcal{X}}}
\newcommand{\hi}{{\hat{\Phi}}}
\newcommand{\data}{{\text{data}}}

\newcounter{rcnt}[section]
\renewcommand{\thercnt}{(\roman{rcnt})}

\def\qt#1{\qquad\text{#1}}

\def\argmin{\mathop{\rm argmin}}
\def\argmax{\mathop{\rm argmax}}


\setlength{\parskip}{1.4 \medskipamount}
%\sloppy
%\linespread{1.3}

\begin{document}

\title{STAT 238 - Bayesian Statistics  \\
  Lecture Six} 
\subtitle{Spring 2026, UC Berkeley}

\author{Aditya Guntuboyina}


\date{02 Feb 2026}

\maketitle

\section{Interpretation of Probability}
Because Bayesian statistics is simply probability theory applied to
inference, understanding the meaning and interpretation of probability
is essential. 

There are broadly two ways of understanding probability: frequentist
and Bayesian. 

\section{Frequentist/Objective Understanding of Probability}
From the frequentist viewpoint, probability is applicable only in the
context of ``random experiments'' (such as tossing coins and rolling
dice). The probability $\P(A)$ of  an event $A$ is defined as the
relative frequency that $A$ occurs in $N$ repeated trials of the
experiment in the limit as $N \rightarrow \infty$:
\begin{equation*}
  \P(A) := \lim_{N \rightarrow \infty} \frac{N_A}{N}
\end{equation*}
where $N_A$ is the number of trials out of $N$ where $A$ occurs. This
definition of probability is the basis of frequentist statistics.

Here are some examples:
\begin{enumerate}
\item The statement $\P(H) = 0.5$ means that the proportion of
  heads in a large number of tosses of the coin approaches $0.5$.
\item The statement
  \begin{equation*}
    \epsilon_1, \dots, \epsilon_n \overset{\text{i.i.d}}{\sim} N(0,
    \sigma^2)
  \end{equation*}
  means that if the experiment generating $\epsilon_1, \dots,
  \epsilon_n$ is repeated a large number of times, the proportion of
  times the values of $(\epsilon_1, \dots, \epsilon_n)$ lie in a set
  $A$ approaches
  \begin{equation*}
    \int_A \prod_{i=1}^n \frac{1}{\sqrt{2 \pi} \sigma} \exp
    \left(-\frac{x_i^2}{2 \sigma^2} \right) dx_1 \dots dx_n
  \end{equation*}
  and this should be true for all subsets $A$ of $\R^n$.
\end{enumerate}

The following are some obvious problems with the frequentist
definition:
\begin{enumerate}
\item It is very restrictive and hardly ever applicable. In many simple
  situations where we would like to use probability, the frequency
  definition is simply does not apply: 
  \begin{enumerate}
  \item Is the suspect X guilty?
  \item What is the chance of rain in Berkeley today?
  \item What is the chance that Y is cancer positive given that
    they tested positive? 
  \end{enumerate}
 For an interesting anecdote about how this restrictive notion does not simply
 make sense in some important problems, see \citet[pages
 43-44]{DeGrootBlackwell}.  
\item Even in situations where the frequency definition is
  seemingly applicable, closer thought might reveal some issues. For
  example, the frequentist probability that a coin comes up
  heads is 0.6 means that 60\% of a large number of tosses of the coin
  should result in $0.6$. But the mechanics of no two tosses are really identical and
  if two tosses are done exactly identically, then we would expect the
  same outcome by the laws of physics. So the term ``identical and
  independent repetitions of an experiment'' is ambiguous.
\end{enumerate}

In the frequentist definition, probability is considered an intrinsic
property of the object under investigation which is only accessible by
an experiment generating samples of infinite size. Thus frequentist
probability is also referred to as ``objective probability''. The 
implication is that we cannot assign it arbitrarily because any
probability assignment that does not agree with the frequency in
infinite trials is wrong. Unfortunately, the actual frequentist
probabililty is seldom known because one cannot generally observe a
large number of repetitions of an experiment and so almost all
probability assignments are wrong from the frequentist point of
view. This is one way of understanding the statistics aphorism: ``All models are
wrong'' (usually attributed to George Box; see
\url{https://en.wikipedia.org/wiki/All_models_are_wrong}).

Here are some quotes by famous statisticians/probabilists illustrating
how widespread frequentist thinking in probability is:

\textit{The numbers $p_r$ should, in fact, be regarded as physical constants of the particular
  die that we are using, and the question as to their numerical values
  cannot be answered by the axioms of probability, any more than the
  size and the weight of the die are determined by the geometrical and
  mechanical axioms. However, experience shows that in a well-made die
  the frequency of any event $r$ in any long series of throws usually
  approaches $1/6$, and accordingly we shall often assume that all the
  $p_r$ are equal to $1/6$}... -- Cram{\'e}r.

  Here is Jaynes's response to the above quote (from page 317 of his
  book): \textit{To a physicist, this statement seems to show utter contempt
  for the known laws of mechanics. The results of tossing a die many
  times do \textbf{not} tell us any definite number characteristic
  only of the die. They tell us also something about how the die was
  tossed. If you toss 'loaded' dice in different ways, you can easily
  alter the relative frequencies of the faces. With only slightly more
  difficulty, you can still do this if your dice are perfectly
  'honest'.}

  Here is a quote by Feller (see page 322 of the Jaynes book)
  illustrating the thinking that bridge hands possess physical
  probabilities and that the uniform probability assignment is a
  convention whose correctness can only be verified by observed
  frequencies in a random experiment : \textit{The number of possible
    distributions of cards in bridge is almost $10^{30}$. Usually we
    agree to consider them as equally probable. For a check of this
    convention more than $10^{30}$ experiments would be required -- a
    billion of billion of years if every living person played one game
    every second, day and night.} -- Feller. 

In spite of these objections, one positive aspect of the frequentist
meaning of probability is that the Rules of Probability follow easily
from this definition. Recall that the rules of probability are: 
\begin{enumerate}
\item $\P(A)$ always  lies between 0 and 1. The probability of an
  impossible event is 0 and the probability of a certain event is 1.
\item Product rule: $\P(A \cap B) = \P(A) \P(B|A) = \P(B) \P(A|B)$.
\item Sum rule: $\P(A \cup B) = \P(A) + \P(B)$ for disjoint events $A$
  and $B$. 
\end{enumerate}


\section{Subjective or Bayesian Understanding of Probability}

In Bayesian statistics, probability is considered a general method of
reasoning under uncertainty. It is applicable to all situations involving 
uncertainty, and is \textit{not} restricted 
to situations involving ``random experiments''.

Further, probability is assumed to have nothing to do with
frequency. This means there is no right or wrong probability
model. Different analysts are welcome to use different models, and
they can be assessed in terms of performance. One can also assess
different models in a Bayesian model selection framework.

Meaning can be assigned to probability statements without relying on
any connection to long-run frequencies. For example, suppose a doctor
assigns a probability of 0.02 to a patient having cancer based on some
background information. This can be quantified as (see \citet[Chapter
3]{Lindley}): \textit{The doctor’s degree of belief
in the uncertain event that the patient has cancer is the same as
the degree of belief of the uncertain event of drawing a red ball
from an urn containing 2 red balls and 98 green balls.}

Lindley goes on to clearly say that there is no frequency
interpretation here. \textit{There is no repetition in this
  definition. The ball is to be 
taken once, and once only, and the long-run frequency of red balls in
repeated drawings is irrelevant. After its withdrawal, the urn and its
contents can go up in smoke for all that it matters.}

For a concrete example, consider the probability assignment $\P\{H
\mid I\} = 0.5$  where $H$ denotes heads (in the context of tossing a
coin) and $I$ denotes some background information. In the Bayesian
context, this is not an informative statement on some long-run
frequency of heads while tossing the coin. Instead, it is an
assignment of probability by some individual for the next coin toss,
based on some information. Since $\P\{H \mid I\} = \P\{T \mid I\}$, it
implies that the background information is symmetric between H and
T. The actual information itself might vary, for example, consider the
following two kinds of background information both of which can
justify this assignment $\P\{H \mid I\} = 0.5$:
\begin{enumerate}
\item \textbf{Information $I_1$}: We don't know anything at all about the
  coin. We don't even know if it really has two sides $H$ and $T$ or
  both of its sides are of only one kind (either $H$ or $T$). In
  addition, we don't know how exactly it will be tossed. 
\item \textbf{Information $I_2$}: We know that it is a ``regular'' coin
  and that it has two sides $H$ and $T$, and that it will be tossed in
  the ``usual'' way. 
\end{enumerate}
In the first case above, it might very well happen that the coin has
both sides H, in which case repeated tossing of the coin will lead to
HHHHH.. so the long run frequency of heads will be 1 (and not
0.5). But this does not invalidate the assignment $\P(H \mid I_1) =
0.5$ because it is not a statement on long run frequency of heads, but
it claims something about the next toss. 

Now consider the third kind of information: We know that this coin has
been tossed a large number of times in the past and it landed heads
70\% of the time. Now is the assumption $\P\{H \mid I\} =
0.5$ justifiable? In this case, the probability we need to calculate
is:
\begin{align}\label{info3}
  \P \left\{X_{n+1} = 1 \mid \frac{X_1 + \dots + X_n}{n} = 0.7 \right\}
\end{align}
where $X_1, \dots, X_n$ represent the historical tosses and $X_{n+1}$
denotes the outcome of the next toss (1 here represents $H$ and 0
represents $T$). The probability \eqref{info3} cannot be assigned
arbitrarily but instead it should be calculated based on some joint
model for $X_1, \dots, X_{n+1}$. Consider the following two models:
\begin{align}\label{mod1}
  X_1, \dots, X_{n+1} \overset{\text{i.i.d}}{\sim} \text{Bernoulli}(0.5)
\end{align}
and
\begin{align}\label{mod1}
  X_1, \dots, X_{n+1} \mid \theta \overset{\text{i.i.d}}{\sim}
  \text{Bernoulli}(\theta) ~~ \text{ and } ~~ \theta \sim
  \text{uniform}(0, 1). 
\end{align}
For the first model, $X_{n+1}$ is independent of $X_1, \dots, X_n$ and
we indeed get $\P(X_{n+1} = 1) = 0.5$. For the second model, the
situation is more interesting, and we will get an answer to
\eqref{info3} that is close to $0.7$ when $n$ is large. So here,
probability still does not have anything to do with frequency, but in
this case, the right kind of model will lead to the frequency
assignment. 

This flexibility with modeling is a positive
feature of the Bayesian framework. However, there remains the question
of justification of the 
rules of probability. If analysts are allowed to come up arbitrary
probability models, then why do they have to carry out
calculations in accordance with the rules of probability. What is the
justification for using the rules of probability? This question was
raised by many people including R. A. Fisher, who in \cite{Fisher1934} writes:  

\textit{Keynes establishes the laws of addition and multiplication of
  probabilities, by stating these laws in the form of definitions of
  the processes of addition and multiplication. The important step of
  showing that, when these probabilities have numerical values,
  ``addition'' and ``multiplication'' are so defined, are equivalent
  to the arithmetical processes ordinarily known by these names, is
  omitted. The omission is an interesting one, since it shows the
  difficulty of establishing the laws of mathematical probability,
  without basing the notion of probability on the concept of
  frequency, for which these laws are really true, and from which they
  were originally derived.}

It turns out that the rules of probability can be justified without
using any connections between probability and frequency. The following
arguments are due to the physicist R. T. Cox and are described in Chapters 1 and
Chapter 2 of \citet{Jaynes}. I will give a sketch of the argument
skipping some important technical details. For the full argument,
please read \citet[Chapter 1 and 2]{Jaynes}.

\section{Justification of the Rules of Probability without using any
  connections to Long-run Frequencies}
As just mentioned, the following argument is due to the physicist
R. T. Cox, and can be read in the book \citet{rtcox}. I will follow
the treatment given in \citet[Chapter 1 and 2]{Jaynes}. 


Let us first remove all restrictions on probabilities and even allow
them to take values outside the interval $[0, 1]$. To avoid confusion,
let us use the term ``plausibilities''. We are assigning
plausibilities of various events (or propositions) conditional on
other events. Let us denote the plausibility of event $A$ conditional
on event $B$ by $(A|B)$. Let us first make the assumption that
plausibilities take values in the set of real numbers (no restriction
now to be in the interval $[0, 1]$) and that a higher value of
plausibility represents a greater belief.

\subsection{Product Rule}
Let us first investigate why the product rule should be true. The
product rule in terms of probabilities states that
\begin{equation*}
  \P(AB|C) = \P(B|C) \P(A|BC)
\end{equation*}
Here $AB$ denotes the event $A \cap B$. Should our plausibilities
satisfy a similar inequality? Let us first assume that the
plausibility $(AB|C)$ should really be determined  by the two
plausibilities $(B|C)$ and $(A|BC)$. This is basically because the
process of deciding that $AB$  is true can be broken down into first
deciding whether $B$ is true and then, having accepted $B$ as true,
deciding whether $A$ is true. We shall therefore assume that there
should be a function $F$ such that
\begin{equation*}
  (AB|C) = F((B|C), (A|BC)). 
\end{equation*}
We also assume that we should use the same function $F$ for all
possible events $A, B, C$ (i.e., we are not using one function $F$ for
some $A, B, C$ while calculating $(AB|C)$ from $(B|C)$ and $(A|BC)$
and using another function $F$ for different $A, B, C$). This means in
particular that
\begin{equation*}
  (AB|C) = F((A|C), (B|AC)). 
\end{equation*}
It is also reasonable to assume that $F(x, y)$ is monotone increasing
in each of its arguments and that it is continuous. If it is not
continuous, then a small change in $(B|C)$ (or $(A|BC)$) might lead to
a large change in $(AB|C)$ which is undesirable.

Now if we have four events $A, B, C, D$, we can write
\begin{equation*}
  (ABC|D) = F((BC|D), (A|BCD)) = F(F((C|D), (B|CD)), (A|BCD)). 
\end{equation*}
We can also write
\begin{equation*}
  (ABC|D) = F((C|D), (AB|CD)) = F((C|D), F((B|CD), (A|BCD))). 
\end{equation*}
We shall now make the following important consistency assumption:
\textbf{If a plausibility can be calculated via two different methods,
then both methods  should give the same answer}. Clearly if this
assumption were violated, then our answer to a plausibility
calculation would depend on the specific method chosen to calculate
and this would be highly undesirable. This assumption immediately
implies that
\begin{equation*}
  F(F((C|D), (B|CD)), (A|BCD)) = F((C|D), F((B|CD), (A|BCD))). 
\end{equation*}
for all $A, B, C, D$. If the individual plausibilities are arbitrary,
we would get the following condition that the function $F$ should
satisfy
\begin{equation*}
  F(F(x, y), z) = F(x, F(y, z)) \qt{for all real numbers $x, y, z$}. 
\end{equation*}
It now turns out the only functions $F$ which satisfy the above
equation are of the form
\begin{equation*}
  F(x, y) = w^{-1}(w(x) w(y))
\end{equation*}
for a positive continuous increasing function $w$. I will skip this
derivation (see Section 2.1, Chapter 2 of \citet{Jaynes}). We thus have
\begin{equation}\label{abc}
  (AB|C) = F((B|C), (A|BC)) = w^{-1} \left(w(B|C) w(A|BC) \right). 
\end{equation}
This is equivalent to
\begin{equation*}
  w(AB|C) = w(B|C) w(A|BC). 
\end{equation*}
Now if we take $B = A$, we get
\begin{equation*}
  w(A|C) = w(A|C) w(A|AC)
\end{equation*}
The event $A|AC$ can be seen as certainty so we get
\begin{equation*}
  w(A|C) = w(A|C) w(\text{certainty})
\end{equation*}
for all $A$ and $C$. This can happen only if
\begin{equation}\label{cert}
  w(\text{certainty}) = 1. 
\end{equation}
Also if we take $B = A^c$ in \eqref{abc}, we get
\begin{equation*}
  w(AA^c | C) = w(A^c|C) w(A|A^c C)
\end{equation*}
$AA^c|C$ and $A|A^cC$ can both be taken to represent impossibility so
we get
\begin{equation*}
  w(\text{impossible}) = w(A^c|C) w(\text{impossible})
\end{equation*}
for all $A$ and $C$ which gives
\begin{equation}\label{impo}
  w(\text{impossible}) = 0. 
\end{equation}
\eqref{cert} and \eqref{impo}, along with the monotonicity of $w$,
imply
\begin{equation}\label{w01}
0 \leq  w(A|B) \leq 1 \qt{for all $A$ and $B$}. 
\end{equation}
We have thus proved that $w(A|B)$ lies always between 0 and 1 (is 0
for impossibility and 1 for certainty) and it satisfies the product
rule of probability: 
\begin{equation}\label{w1}
  w(AB|C) = w(A|C) w(B|AC) = w(B|C) w(A | BC). 
\end{equation}
In other words, if we apply this this function $w$ to our plausibilities, then the
resulting assignments satisfy the first two rules of
probability.

\subsection{Sum Rule}

Next the goal is to derive the sum rule. We will first derive the sum rule
in the simplified form: $\P(A^c) = 1 - \P(A)$, and then prove it in
the general case. We shall discuss this argument in the next lecture. 



\begin{thebibliography}{99}

\bibitem[{Cox(2001)}]{rtcox}Cox, R. T. (2001)
       \textit{Algebra of Probable Inference}. John Hopkins University
       Press, 2001. 

\bibitem[Fisher(1934)]{Fisher1934}
R.~A. Fisher,
\newblock ``Probability, likelihood and quantity of information in the logic of
uncertain inference,''
\newblock \emph{Proceedings of the Royal Society of London. Series~A,
Containing Papers of a Mathematical and Physical Character}, vol.~146,
no.~856, pp.~1--8, 1934.

     \bibitem[Jaynes(2003)]{Jaynes} Jaynes, E. T, (2003)
  \textit{Probability theory: the logic of science}. Cambridge
  University Press, 2003.

%  \bibitem[Sinai(1992)]{Sinai} Sinai, Y. G, (1992)
%  \textit{Probability theory: an introductory course}. Springer
%  Verlag, 1992.  

\bibitem[{deGroot(1986)}]{DeGrootBlackwell}DeGroot, M. H. (1986)
       A conversation with David Blackwell.
\newblock \emph{Statistical Science}, \textbf{1}, 40--53.

%         \bibitem[Mosteller(1987)]{Mosteller} Mosteller, F, (1987)
%  \textit{Fifty challenging problems in probability with
%    solutions}. Courier Corporation, 1987.

%    \bibitem[MacKay(2003)]{MacKay} MacKay, D. J. C., (2003)
%  \textit{Information theory, inference, and learning
%  algorithms}. Cambridge University Press, 2003.

\bibitem[{Lindley(2013)}]{Lindley}Lindley, D. V. (2013)
       \textit{Understanding Uncertainty}. John Wiley and sons, 2013.


  \end{thebibliography}







\end{document}

%%% Local Variables:
%%% mode: latex
%%% TeX-master: t
%%% End:
